% TOP_PROBLEM: {"id": 30006770, "title": "Smooth anisotropic Calderón uniqueness from full boundary data", "category_id": 9, "category": {"id": 9, "name": "pde", "display_name": "Partial Differential Equations", "description": "PDEs and their applications in physics and geometry.", "slug": "pde", "order_index": 9, "created_at": "2026-07-31T15:26:25.671Z"}, "status": "open"}
% FIRST_POSTED: 2026-10-06
% ATTEMPT_STATUS: solved
% OPENAI_CLAIM: {"schema_version":1,"source_repo":"https://github.com/openai/math","source_commit":"adc7f1241b42e322a6451854ab7e4b4c146bf78a","release_date":"2026-10-06","match":"stronger","resolution":"proved","independently_reviewed":false,"second_pass":"agreed","families":[{"family":"365","title":"Joint metric and connection recovery from one boundary patch","subject":"Partial differential equations","manuscripts":[{"title":"Determination of a metric and a unitary connection from one boundary patch","pdf_path":"preprints/Determination-of-a-metric-and-a-unitary-connection-from-one-boundary-patch-October-5-2026/paper.pdf","readme_path":"preprints/Determination-of-a-metric-and-a-unitary-connection-from-one-boundary-patch-October-5-2026/README.md","date":"2026-10-05","theorem_ref":"thm:main (sections/01-introduction.tex:39)","statement_sources":[{"path":"preprints/Determination-of-a-metric-and-a-unitary-connection-from-one-boundary-patch-October-5-2026/build/sections/01-introduction.tex","line":39,"label":"thm:main"}]},{"title":"Smooth anisotropic uniqueness in the Calderón problem from one boundary patch","pdf_path":"preprints/Smooth-Anisotropic-Uniqueness-in-the-Calderon-Problem-from-One-Boundary-Patch-September-24-2026/paper.pdf","readme_path":"preprints/Smooth-Anisotropic-Uniqueness-in-the-Calderon-Problem-from-One-Boundary-Patch-September-24-2026/README.md","date":"2026-09-24","theorem_ref":"main:full (sections/01-introduction.tex:84); main:patch (sections/01-introduction.tex:54)","statement_sources":[{"path":"preprints/Smooth-Anisotropic-Uniqueness-in-the-Calderon-Problem-from-One-Boundary-Patch-September-24-2026/build/sections/01-introduction.tex","line":84,"label":"main:full"},{"path":"preprints/Smooth-Anisotropic-Uniqueness-in-the-Calderon-Problem-from-One-Boundary-Patch-September-24-2026/build/sections/01-introduction.tex","line":54,"label":"main:patch"}]}],"lean":{"doc_path":"lean/docs/365.md","comparators":[],"declarations":[],"covers_main_theorem":false},"reasoning_trace":null}]}
\documentclass[11pt]{article}
\usepackage[margin=25mm]{geometry}
\usepackage{fontspec}
\setmainfont{Latin Modern Roman}
\usepackage{amsmath,amssymb,mathtools}
\usepackage{unicode-math}
\setmathfont{Latin Modern Math}
\usepackage{enumitem}
\usepackage{xurl,hyperref,bookmark}
\hypersetup{colorlinks=true,urlcolor=blue}
\setcounter{secnumdepth}{0}
\setlength{\parindent}{0pt}
\setlength{\parskip}{5pt}
\setlength{\emergencystretch}{3em}
\begin{document}
\section*{Smooth anisotropic Calderón uniqueness from full boundary data}

\subsection{OpenAI's claim}
\textit{Recovery of a metric and connection}: Zero-frequency measurements with input and observation on one arbitrary open boundary patch are claimed to recover both a smooth metric and a unitary rank-two connection; recovery is up to a diffeomorphism and unitary gauge that fix the measured patch.

\textit{One-patch anisotropic uniqueness}: The paper claims determination of a smooth Riemannian metric in dimension at least three from its zero-frequency energy form on any one open boundary patch; the only ambiguity is a diffeomorphism equal to the identity on that patch.

\subsection{Mathematical statement of the claimed result}
\paragraph{Mathematical claim 1.} Fix $n\ge3$ and a smooth connected manifold $M$ of that dimension, assumed compact with $\partial M$ smooth and nonempty. Choose a relatively open $\Gamma\subset\partial M$ which is nonempty and proper. On the trivial Hermitian rank-two bundle, set $d_A=d+A$, $A\in C^\infty(T^*M\otimes\mathfrak u(2))$, and define the local Dirichlet energy form using solutions of $d_A^{*g}d_Au=0$ with boundary data supported in $\Gamma$. If two smooth metrics and connections give the same form for every $f,h\in C_c^\infty(\Gamma;\mathbb C^2)$, the claim supplies smooth $\Phi:M\to M$ and $U:M\to U(2)$ with \[\Phi|_\Gamma=\operatorname{id},\quad U|_\Gamma=I,\quad g_2=\Phi^*g_1,\quad A_2=U^{-1}(\Phi^*A_1)U+U^{-1}dU.\] Thus the ambiguity in recovering the two geometric unknowns consists of the displayed diffeomorphism and gauge, each trivial on the observation patch.

\paragraph{Mathematical claim 2.} Fix a dimension $n\ge3$. Assume the smooth $n$-manifold $M$ is compact and connected, with smooth nonempty boundary. For a proper relatively open observation patch $\Gamma\ne\varnothing$, the claim asserts \[\Lambda_{g_1,\Gamma}=\Lambda_{g_2,\Gamma}\quad\Longrightarrow\quad g_2=\Phi^*g_1,\qquad\Phi|_\Gamma=\operatorname{id},\] for some diffeomorphism smooth through the boundary. Both metrics are smooth, and the local maps are the harmonic Dirichlet energy forms for data supported in $\Gamma$. Equality for all full-boundary data gives a diffeomorphism fixing the entire boundary. This concerns smooth anisotropic metrics, not arbitrary bounded measurable scalar conductivities.

\subsection{Outline of the argument}
\textit{Recovery of a metric and connection}: Boundary factorization identifies normal-gauge jets, and a common exterior cap starts equality of Green matrices. Source solutions provide separating points, harmonic jets, and bundle frames. Mixed-kernel continuation in the product, followed by moving-sphere transfer and angular coercivity, extends the maximal local metric-and-bundle identification throughout the manifold.

\textit{One-patch anisotropic uniqueness}: Boundary jets and a common exterior cap identify actual Green kernels in an initial region. Remote source solutions produce harmonic coordinates and enough second-order jets. A mixed Green-kernel continuation theorem extends metric equality across a smooth frontier, so the maximal paired coordinate region cannot stop in the interior.

\subsection{Scope relative to this problem}
\textit{Recovery of a metric and connection}: With zero connections the vector-valued patch theorem determines the scalar metric up to a diffeomorphism fixing that patch. The site full-boundary gauge must fix the entire boundary; unlike the scalar companion paper, this manuscript main theorem does not explicitly provide that all-boundary marking conclusion.

\textit{One-patch anisotropic uniqueness}: Corollary main:full gives exact full-boundary smooth anisotropic uniqueness with a diffeomorphism fixing the entire boundary in every n>=3, matching the site gauge and manifold hypotheses. The main theorem further treats one arbitrary nonempty proper open patch, with the corresponding patch gauge.

The claimed result implies the stated question and also supplies the additional strength described above.

\subsection{Formal verification}
Family 365 describes Lean coverage for the measurable scalar-conductivity obstruction; smooth anisotropic and coupled connection recovery are separate.

\textit{Recovery of a metric and connection}: No manuscript-specific Comparator is listed in the pinned inventory.

The documented selection does not cover the entire manuscript assertion relevant here.

\textit{One-patch anisotropic uniqueness}: No manuscript-specific Comparator is listed in the pinned inventory.

The documented selection does not cover the entire manuscript assertion relevant here.

This site has not run the Lean development. These links describe OpenAI’s selected formal statements and should be read with their stated scope.

\subsection{Status caveat}
OpenAI reports that an internal model produced the manuscripts. Their arguments have not been independently verified as proofs by this site. Any portion outside the documented formal scope may contain mistakes, and the catalogue label reports the announced mathematical claim.

\subsection{Sources}
\begin{itemize}
\item OpenAI, \emph{Determination of a metric and a unitary connection from one boundary patch}, 2026-10-05; citation key \texttt{OAI:Determination-of-a-metric-and-a-unitary-connection-from-one-boundary-patch-October-5-2026}. \href{https://github.com/openai/math/blob/adc7f1241b42e322a6451854ab7e4b4c146bf78a/preprints/Determination-of-a-metric-and-a-unitary-connection-from-one-boundary-patch-October-5-2026/paper.pdf}{PDF (GitHub)}; \href{https://github.com/openai/math/raw/adc7f1241b42e322a6451854ab7e4b4c146bf78a/preprints/Determination-of-a-metric-and-a-unitary-connection-from-one-boundary-patch-October-5-2026/paper.pdf}{PDF (direct)}; \href{https://github.com/openai/math/blob/adc7f1241b42e322a6451854ab7e4b4c146bf78a/preprints/Determination-of-a-metric-and-a-unitary-connection-from-one-boundary-patch-October-5-2026/README.md}{Manuscript page and citation}; \href{https://github.com/openai/math/blob/main/preprints/Determination-of-a-metric-and-a-unitary-connection-from-one-boundary-patch-October-5-2026/paper.pdf}{Latest version}; Statement source: \href{https://github.com/openai/math/blob/adc7f1241b42e322a6451854ab7e4b4c146bf78a/preprints/Determination-of-a-metric-and-a-unitary-connection-from-one-boundary-patch-October-5-2026/build/sections/01-introduction.tex\#L39}{thm:main}.
\item OpenAI, \emph{Smooth anisotropic uniqueness in the Calderón problem from one boundary patch}, 2026-09-24; citation key \texttt{OAI:Smooth-Anisotropic-Uniqueness-in-the-Calderon-Problem-from-One-Boundary-Patch-September-24-2026}. \href{https://github.com/openai/math/blob/adc7f1241b42e322a6451854ab7e4b4c146bf78a/preprints/Smooth-Anisotropic-Uniqueness-in-the-Calderon-Problem-from-One-Boundary-Patch-September-24-2026/paper.pdf}{PDF (GitHub)}; \href{https://github.com/openai/math/raw/adc7f1241b42e322a6451854ab7e4b4c146bf78a/preprints/Smooth-Anisotropic-Uniqueness-in-the-Calderon-Problem-from-One-Boundary-Patch-September-24-2026/paper.pdf}{PDF (direct)}; \href{https://github.com/openai/math/blob/adc7f1241b42e322a6451854ab7e4b4c146bf78a/preprints/Smooth-Anisotropic-Uniqueness-in-the-Calderon-Problem-from-One-Boundary-Patch-September-24-2026/README.md}{Manuscript page and citation}; \href{https://github.com/openai/math/blob/main/preprints/Smooth-Anisotropic-Uniqueness-in-the-Calderon-Problem-from-One-Boundary-Patch-September-24-2026/paper.pdf}{Latest version}; Statement source: \href{https://github.com/openai/math/blob/adc7f1241b42e322a6451854ab7e4b4c146bf78a/preprints/Smooth-Anisotropic-Uniqueness-in-the-Calderon-Problem-from-One-Boundary-Patch-September-24-2026/build/sections/01-introduction.tex\#L84}{main:full}; Statement source: \href{https://github.com/openai/math/blob/adc7f1241b42e322a6451854ab7e4b4c146bf78a/preprints/Smooth-Anisotropic-Uniqueness-in-the-Calderon-Problem-from-One-Boundary-Patch-September-24-2026/build/sections/01-introduction.tex\#L54}{main:patch}.
\item Family 365: \href{https://github.com/openai/math/blob/adc7f1241b42e322a6451854ab7e4b4c146bf78a/lean/docs/365.md}{Lean scope}.
\item \href{https://github.com/openai/math/blob/adc7f1241b42e322a6451854ab7e4b4c146bf78a/CONTENTS.md}{OpenAI catalogue entry}.
\end{itemize}
\end{document}
